\documentclass[12pt,a4paper,titlepage]{article}
\usepackage[utf8]{inputenc}
\usepackage[T1]{fontenc}
\usepackage{lmodern}
\usepackage{setspace}
\usepackage{geometry}
\usepackage{amsmath,amssymb}
\usepackage{longtable}
\usepackage{hyperref}
\usepackage{enumitem}
\usepackage{natbib}
\usepackage{multirow}
\usepackage{graphicx} % Required for \resizebox
\usepackage{booktabs} % Required for \toprule, \midrule, \bottomrule
\usepackage{listings}
\usepackage{xcolor}
\lstdefinelanguage{TypeScript}{
  keywords={abstract, any, as, boolean, break, case, catch, class, console, const, continue, debugger, declare, default, delete, do, else, enum, export, extends, false, finally, for, from, function, get, if, implements, import, in, instanceof, interface, let, module, namespace, never, new, null, number, object, package, private, protected, public, readonly, require, return, set, static, string, super, switch, this, throw, true, try, type, typeof, var, void, while, with, yield, async, await, Promise},
  morecomment=[l]{//},
  morecomment=[s]{/*}{*/},
  morestring=[b]',
  morestring=[b]",
  morestring=[b]`,
  ndkeywords={class, export, boolean, throw, implements, import, this},
  keywordstyle=\color{blue}\bfseries,
  ndkeywordstyle=\color{darkgray}\bfseries,
  identifierstyle=\color{black},
  sensitive=false,
  commentstyle=\color{gray}\ttfamily,
  stringstyle=\color{teal}\ttfamily,
}

\lstset{
    language=Python,
    basicstyle=\ttfamily\small,
    keywordstyle=\color{blue}\bfseries,
    commentstyle=\color{gray}\itshape,
    stringstyle=\color{teal},
    showstringspaces=false,
    numbers=left,
    numberstyle=\tiny\color{gray},
    stepnumber=1,
    numbersep=5pt,
    backgroundcolor=\color{gray!5},
    frame=single,
    framerule=0.5pt,
    tabsize=4,
    breaklines=true,
    breakatwhitespace=true,
    captionpos=b,
    morekeywords={self, True, False, None}
}

\geometry{margin=1in}
\onehalfspacing

\title{LLM-Orchestrated Neighbor-Enriched Forecasting for NIFTY Equity Returns}

\author{
    Bijit Mondal\\
    % \texttt{hello@bijit.xyz}
}
\date{} % Explicitly leaving this empty removes the date from the title page

\begin{document}

\maketitle

\begin{abstract}
Cross-stock return forecasting is central to quantitative research in equity markets, yet practitioners and researchers face a persistent bottleneck: manually discovering which other stocks carry predictive signal for a target ticker is labor-intensive, error-prone, and rarely reproducible. Many existing approaches fix a handful of neighbors or ignore calendar alignment entirely, leading to invalid comparisons and fragile conclusions. In this paper, we introduce ECHO (Equity Cross-stock Heuristic Orchestrator), the first LLM-driven agentic experimentation platform for systematic cross-stock neighbor discovery and multi-horizon cumulative return forecasting on the Indian NIFTY equity universe. The framework comprises an Orchestrator that proposes diverse neighbor sets using metadata, sector graphs, and memory of past experiments; a Forecast Service that trains a fixed XGBoost regressor on calendar-aligned daily returns with explicit chronological splits and naive baselines (zero-return and train-mean); an Experiment Logger that records every proposal with full traceability; and an Analyzer that aggregates results into a calibrated dashboard for honest evaluation. Rather than claiming a superior model, ECHO treats the forecaster as a rigorous measuring instrument and focuses on automating the full research loop: propose $\rightarrow$ measure $\rightarrow$ remember $\rightarrow$ refine. Across thousands of logged experiments on NIFTY constituents, the platform demonstrates clear process-level superiority over manual and fixed-neighbor baselines while transparently reporting when even the best configurations fail to beat the zero-return baseline—a common and scientifically valid outcome in noisy return data. By turning forecasting into a repeatable, auditable experimentation pipeline, ECHO bridges the gap between manual workflows and production-grade quantitative research.

\vspace{0.5em}
% \textbf{GitHub:} \url{https://github.com/BMondal-dev/financial-agent} \quad

% \textbf{Corresponding Author:} hello@bijit.xyz
\end{abstract}

\clearpage
\tableofcontents
\clearpage

\section{Introduction}

Return forecasting in equity markets guides investment decisions, risk management, and portfolio construction across hedge funds, asset managers, and academic research. In practice, researchers routinely analyze thousands of short-to-medium daily return series that exhibit non-stationarity, regime shifts, and cross-sectional dependencies. The dominant cost lies not in fitting a single model, but in the labor-intensive process of hypothesis generation (which stocks might help predict a target?), rigorous evaluation (proper calendar alignment, chronological splits, and strong baselines), and institutional memory (what worked before and why?). Many existing approaches and baseline pipelines shortcut this workflow—often aligning series by row index instead of trading date, omitting naive baselines, and reporting a single MAE number without traceability—leading to invalid comparisons and fragile conclusions.

Most advances in return forecasting arrive either as expert models tuned to a fixed feature set or as universal deep-learning forecasters that still leave the broader pipeline of feature (neighbor) discovery and evaluation untouched. AutoML-style approaches focus narrowly on hyperparameter search while ignoring the richer hypothesis space of cross-stock predictors. Meanwhile, naive LLM-based forecasters treat time series as language but lack grounding in market calendars and explicit baselines, making their outputs difficult to trust or reproduce. These limitations motivate an agentic approach—one that treats cross-stock return forecasting as a sequential decision process over neighbor proposal, rigorous scoring, logging, and refinement, with explicit planning, tool use, and transparent rationales.

To this end, we introduce ECHO (Equity Cross-stock Heuristic Orchestrator), the first end-to-end, agentic experimentation platform that automates the entire workflow a human quantitative researcher would follow for discovering predictive cross-stock neighbors on the NIFTY universe. Rather than committing to a single universal model, ECHO orchestrates four specialized components that work in concert. First, a metadata engine builds a transparent feature layer per stock by combining Yahoo Finance fundamentals (via \texttt{ticker.info}) with computed market signals (volatility, momentum, drawdown, trend regime, static and rolling correlations), then compiles similarity sets and a weighted stock graph. The Orchestrator (powered by a structured-output LLM) reasons over this structured context, graph neighborhoods, and memory of past outcomes to propose candidate neighbor sets. Furthermore, by integrating qualitative market news sentiment and broader sector narratives, the Orchestrator evaluates the macroeconomic mood alongside statistical linkages. This combined quantitative-qualitative reasoning, augmented by memory of past experiment outcomes, enables the system to generate highly contextualized neighbor hypotheses. Next, the Forecast Service builds a properly calendar-aligned dataset (inner join on trading \texttt{Date}), trains the \emph{same} XGBoost regressor on lagged features and neighbor returns, evaluates on a strict chronological holdout, and computes honest metrics against zero-return and train-mean baselines. The Experiment Logger records every run with full provenance (rationale, neighbor list, horizon, source, and \texttt{mae\_for\_ranking} penalty when the model fails to beat zero). Finally, the Analyzer aggregates the append-only log into a calibrated dashboard that surfaces baseline comparisons (A = target-only features, B = fixed correlation neighbors, Ours = best agent-proposed set), sector influence matrices, predictor frequency, and backtest overlays—turning raw experiment history into actionable intelligence.

Across extensive experiments on NIFTY 50 constituents, ECHO consistently outperforms unstructured manual workflows and fixed-neighbor baselines in terms of hypothesis coverage, evaluation hygiene, and reproducibility. While many configurations do not beat the zero-return baseline (a scientifically honest outcome in noisy daily return data), the platform’s Baseline A/B/Ours comparisons and calibrated leaderboards reveal where agent-chosen neighbors provide measurable lift over naive alternatives. The Baseline A vs B vs Ours comparison (target-only vs fixed correlation vs agent-proposed) shows that target-only features capture much of the autoregressive structure, but the agent-proposed neighbors consistently outperform static correlation neighbors, demonstrating that smart neighbor selection avoids poor choices and adds interpretable cross-stock signal. Ablations confirm that calendar alignment, explicit baselines, and multi-round memory each contribute materially to the quality of discovered configurations. The generated dashboard further provides transparent, visual reports that make the entire experimentation process interpretable and extensible.

From a research methodology perspective, the central challenge is combinatorial. For a universe of \(N\) stocks and neighbor set size \(k\), the number of possible cross-stock hypotheses for a single target is \(\binom{N-1}{k}\), which quickly becomes intractable even for moderate \(N\). Human researchers therefore rely on coarse shortcuts (same sector, top historical correlation, or discretionary intuition), but such shortcuts frequently miss transient cross-sector lead-lag effects and regime-dependent dependencies. ECHO is designed to address this search problem directly: rather than collapsing the hypothesis space to one static rule, it performs structured proposal generation guided by metadata, graph priors, and memory, then validates each proposal under identical forecasting conditions.

A second challenge is data integrity under financial market frictions. Indian equities exhibit non-uniform trading calendars across symbols due to listing dates, corporate events, and occasional missing observations; benchmark indices follow their own data availability patterns. Naively joining time series by row order or filling missing points without discipline can induce silent leakage and spurious correlation. Consequently, our framework treats calendar alignment as a first-class constraint and keeps evaluation chronologically strict by construction. This design choice is not merely an implementation detail: it is a scientific guardrail that ensures measured gains are attributable to improved neighbor hypotheses rather than artifacts of misaligned data.

The NIFTY universe offers a particularly suitable testbed for this thesis. It is sector-diverse (financials, energy, information technology, consumer, industrial, healthcare), sufficiently liquid for stable daily observations, and has experienced materially different macro regimes within the study window, including pandemic shock, recovery, inflationary pressure, and policy-rate transitions. These shifts create a realistic environment in which cross-stock relationships evolve over time. Static all-period associations can therefore be informative yet incomplete, motivating the joint use of long-horizon correlation structure and rolling, short-horizon co-movement signals.

Our design also emphasizes semantic clarity between \emph{descriptors} and \emph{derived signals}. Certain fields, such as \texttt{market\_cap} and \texttt{beta}, are externally sourced descriptors from Yahoo Finance and should be interpreted as contextual priors, not model outputs. In contrast, momentum, volatility, drawdown, trend regime, and rolling correlations are computed transformations of observed prices. This separation enables more transparent reasoning by the Orchestrator and more interpretable ablations: researchers can explicitly analyze whether improvements stem from stable fundamental context, dynamic return structure, or their interaction in the weighted similarity graph.

Finally, ECHO adopts a process-level view of progress. In noisy return forecasting, single-run wins can be accidental; what matters is whether the research loop itself is rigorous, repeatable, and cumulative. By fixing the forecaster, enforcing leakage-safe evaluation, and logging every proposal with rationale and outcomes, ECHO turns neighbor discovery into an auditable experimental science workflow. This perspective aligns with practical quantitative research, where trustworthy iteration and institutional memory are often more valuable than isolated point improvements on one backtest.

In the same spirit, we intentionally use a fixed \texttt{XGBRegressor} and a stateless, request-scoped training policy instead of persisting a single global model file. In ECHO, the neighbor set is part of the hypothesis itself, so changing neighbors changes feature identity. For example, a model trained with neighbors \(\{A,B,C\}\) expects features such as \texttt{A\_lag\_1}, \texttt{B\_lag\_1}, and \texttt{C\_lag\_1}; if a later request uses \(\{A,D,E\}\), the resulting columns become \texttt{A\_lag\_1}, \texttt{D\_lag\_1}, and \texttt{E\_lag\_1}, and the saved model is semantically mismatched. Moreover, forecasting always depends on the latest bars, which requires fresh lag and rolling-feature construction even when a model checkpoint exists. We therefore prioritize proposal-specific retraining and methodological consistency, and avoid sequence-heavy alternatives (e.g., LSTM) whose added training complexity can confound whether gains come from better neighbors or from model-class tuning.


\textbf{Our main contributions are as follows:}
\begin{enumerate}
    \item We introduce ECHO (Equity Cross-stock Heuristic Orchestrator), the first end-to-end agentic experimentation platform for cross-stock neighbor discovery and return forecasting on the NIFTY universe.
    \item We formalize a transparent metadata thesis by separating direct Yahoo Finance attributes (e.g., \texttt{market\_cap}, \texttt{beta}) from computed signals (volatility, momentum, relative strength, drawdowns, trend regime, static/rolling correlations, and similarity sets).
    \item We define an explicit weighted stock-similarity graph that combines correlation strength with sector, volatility-bucket, and market-cap-bucket agreement, and use it as a structural prior for neighbor proposal.
    \item We formalize a leakage-safe, calendar-aligned evaluation protocol with explicit naive baselines, cumulative returns over the prediction horizon, and a ranking penalty that prevents celebration of trivial performance.
    \item We justify a fixed-model experimentation strategy (\texttt{XGBRegressor} with stateless retraining) that avoids schema mismatch under changing neighbor sets and keeps predictions tied to current market data.
    \item We demonstrate that LLM-guided neighbor proposals, when rigorously scored by a fixed XGBoost yardstick, consistently outperform fixed-correlation baselines and approach target-only performance on many targets while maintaining full traceability.
\end{enumerate}


\section{Related Work}

\subsection{Equity Return Forecasting and Cross-Stock Prediction}

Equity return forecasting has long been a cornerstone of quantitative finance. Early work focused on classical cross-sectional factors such as size, value, and momentum \citep{fama1993common,jegadeesh1993returns}. More recent studies have leveraged machine learning to exploit the high-dimensional nature of firm characteristics, demonstrating substantial improvements in out-of-sample predictability \citep{gu2020empirical,chen2023deep}. Deep learning models have further advanced the field by capturing non-linear patterns in the cross-section of returns \citep{abe2018deep,zhang2024}.

A growing body of research highlights the importance of cross-stock relationships. Economically linked firms exhibit return predictability due to gradual information diffusion \citep{cohen2008supply,menzly2010market}. Graph-based approaches have emerged to explicitly model these interdependencies, constructing stock correlation or industry networks and feeding them into graph neural networks \citep{matsunaga2019graph,peng2020}. More recent work incorporates LLM-augmented semantic networks to refine economic linkages and explores cross-stock trend integration for enhanced predictive modeling. Despite these advances, most methods still rely on manually curated or statically defined neighbor sets, lack rigorous calendar alignment, and seldom incorporate explicit baselines or institutional memory—leaving a significant gap in automated, reproducible hypothesis discovery.

\subsection{Agentic and LLM-Driven Approaches for Time Series and Financial Forecasting}

Large language models have recently been applied to time series forecasting, either by reprogramming pretrained LLMs \citep{jin2023time, zhou2023} or by treating numerical series as language tokens. While these models show promise in zero-shot settings, they remain primarily model-centric and rarely address the full experimentation pipeline.

The rise of multi-agent systems has opened new possibilities for complex analytical tasks. Frameworks such as CAMEL \citep{li2023camel} and AutoGen \citep{wu2023autogen} enable collaborative LLM agents to tackle reasoning, planning, and tool use. In finance, multi-agent architectures have been explored for trading decision-making and sentiment analysis.

Most relevant to our work is TimeSeriesScientist \citep{zhao2025timeseriesscientist}, the first LLM-driven agentic framework for general time series forecasting (available at \url{https://arxiv.org/abs/2510.01538}). It employs specialized agents for curation, planning, forecasting, and reporting, demonstrating strong performance and interpretability across benchmarks. However, it targets univariate series in non-financial domains and does not focus on cross-stock neighbor discovery or equity-specific challenges such as calendar alignment and baseline-aware ranking.

Recent agentic systems for financial time series emphasize planning and data-centric refinement but stop short of a closed-loop experimentation platform that systematically proposes, scores, logs, and refines cross-stock hypotheses under rigorous evaluation protocols.

ECHO advances this line of research by introducing a full agentic experimentation loop specifically designed for cross-stock return forecasting. Unlike prior work that centers on model improvement, ECHO keeps the underlying forecaster fixed (XGBoost) and automates the harder problem of hypothesis generation and honest evaluation, providing the first end-to-end, reproducible pipeline for neighbor-enriched equity forecasting.


\section{ECHO Framework}

ECHO acts as a human quantitative researcher, systematically performing hypothesis generation, rigorous evaluation, logging, and refinement by leveraging LLM reasoning abilities. ECHO integrates four specialized components that work in concert: (1) Orchestrator proposes diverse neighbor sets using metadata, sector graphs, and memory of past experiments; (2) Forecast Service builds calendar-aligned datasets and evaluates a fixed XGBoost regressor with explicit baselines; (3) Experiment Logger records every run with full provenance; and (4) Analyzer aggregates the append-only log into a calibrated dashboard for transparent insights. This design transforms cross-stock return forecasting into an adaptive, reproducible, and interpretable experimentation pipeline.

\begin{figure}[htbp]
    \centering
    \includegraphics[width=0.95\textwidth]{ECHO_Framework.png}
    \caption{Overview of the ECHO framework. The agentic experimentation loop consists of four components: the Orchestrator proposes neighbor sets using metadata, sector graph, and memory; the Forecast Service performs calendar-aligned evaluation with a fixed XGBoost regressor and explicit baselines; the Experiment Logger records full provenance; and the Analyzer produces a calibrated dashboard. Best configurations are fed back as memory for iterative refinement.}
    \label{fig:echo-overview}
\end{figure}

\subsection{Problem Formulation}

Let \( r_t^{(i)} = \frac{P_t^{(i)} - P_{t-1}^{(i)}}{P_{t-1}^{(i)}} \) denote the daily return of stock \( i \) at time \( t \). Given a target stock \( i \), a prediction horizon \( H \), and a candidate set of neighbor stocks \( \mathcal{N} \), the goal is to forecast the \( H \)-day cumulative forward return \( \hat{R}_{t:t+H}^{(i)} \), defined as the total return from holding the stock for \( H \) trading days:
\[
R_{t:t+H}^{(i)} = \frac{P_{t+H}^{(i)}}{P_t^{(i)}} - 1.
\]
The objective is to minimize the mean absolute error (MAE) on a strict chronological holdout set:

\[
\text{MAE} = \frac{1}{|\mathcal{T}_{\text{test}}|} \sum_{t \in \mathcal{T}_{\text{test}}} \left| R_{t:t+H}^{(i)} - \hat{R}_{t:t+H}^{(i)} \right|
\]

\subsection{Metadata Construction and Similarity Graph Logic}

Before orchestration begins, ECHO creates a metadata object for each stock. For each symbol \(s\), close prices are read from local historical CSV files, while selected fundamentals are fetched directly from Yahoo Finance's \texttt{ticker.info}: \texttt{sector}, \texttt{industry}, \texttt{market\_cap}, and \texttt{beta}. These fields are not predicted by the model; they are external descriptors used for hypothesis generation.

\paragraph{Return-loading protocol.}
For each stock, the service ingests raw CSV records after discarding the three non-data header rows, retaining the pair \((\text{Date},\text{Close})\) as the observed market series. Dates are parsed as calendar timestamps and closes are coerced to real-valued prices; malformed entries are mapped to missing values and excluded. The cleaned table is then indexed by \(\text{Date}\), which enforces a time-ordered, calendar-consistent representation for all downstream joins and lag operations. Daily arithmetic returns are computed as
\[
r_t^{(s)} = \frac{P_t^{(s)}}{P_{t-1}^{(s)}} - 1.
\]
Because the first timestamp has no predecessor \(P_{t-1}^{(s)}\), it is removed by construction, and the function outputs a date-indexed univariate return series \(\{r_t^{(s)}\}\).

\paragraph{Forecast-service endpoint roles.}
The orchestrator-facing endpoints expose complementary views of the search space.

\begin{table}[htbp]
\centering
\small
\caption{FastAPI endpoints used by the forecasting workflow.}
\label{tab:forecast-endpoints}
\begin{tabular}{p{0.25\linewidth} p{0.20\linewidth} p{0.48\linewidth}}
\toprule
Endpoint & Source & Returned object \\
\midrule
\texttt{/metadata/\{symbol\}} & \texttt{metadata.json} & Raw structured descriptor set (e.g., \texttt{same\_sector}, \texttt{same\_market\_cap\_bucket}, \texttt{top\_correlated} with correlation values). \\
\texttt{/candidate-neighbors/\{symbol\}} & \texttt{metadata.json} & Flat de-duplicated candidate list formed by the union of sector peers, market-cap-bucket peers, and top-correlated symbols. \\
\texttt{/graph-neighbors/\{symbol\}} & \texttt{stock\_graph.json} & Top-5 graph neighbors sorted by learned edge weight (structural similarity prior). \\
\texttt{/best-neighbors/\{symbol\}} & \texttt{experiments.json} & Top-5 historically best neighbor configurations for that target (lowest MAE among unique sets, optionally horizon-filtered). \\
\bottomrule
\end{tabular}
\end{table}


Metadata and candidate endpoints define the raw pool; graph neighbors inject structural priors; best-neighbor history provides empirical evidence from prior runs. News sentiment (queried separately from Google News in the orchestration layer) contributes a contemporaneous qualitative context signal.

Let \(P_t^{(s)}\) be close price and \(r_t^{(s)} = \frac{P_t^{(s)} - P_{t-1}^{(s)}}{P_{t-1}^{(s)}}\). The metadata pipeline computes:
\begin{itemize}
    \item \textbf{Volatility (30-day)}: \(\sigma_{30}^{(s)} = \operatorname{std}(r_{t-29:t}^{(s)})\), using the latest rolling window.
    \item \textbf{Momentum}: \(m_k^{(s)} = \frac{P_t^{(s)}}{P_{t-k}^{(s)}} - 1\), for \(k \in \{20, 60\}\).
    \item \textbf{Relative Strength vs NIFTY}: \(RS_t^{(s)} = \frac{\prod_{u \le t}(1 + r_u^{(s)})}{\prod_{u \le t}(1 + r_u^{(\text{NIFTY})})}\), computed on date-aligned returns.
    \item \textbf{Drawdown metrics}: with rolling 252-trading-day high \(M_t^{(s)} = \max_{u \in [t-251, t]} P_u^{(s)}\), drawdown is \(DD_t^{(s)} = \frac{P_t^{(s)} - M_t^{(s)}}{M_t^{(s)}}\). Then \texttt{max\_drawdown\_1y} is \(\min_t DD_t^{(s)}\), and \texttt{current\_drawdown} is the latest \(DD_t^{(s)}\).
    \item \textbf{Trend regime}: using \(MA_{50}\) and \(MA_{200}\), we assign \texttt{Bullish} if \(P_t > MA_{50} > MA_{200}\), \texttt{Bearish} if \(P_t < MA_{50} < MA_{200}\), and \texttt{Sideways} otherwise.
\end{itemize}

Cross-sectional buckets are assigned using \texttt{percentile\_bucket()} over all stocks in the current universe:
\[
\pi_s(x) = \frac{1}{|\mathcal{S}|}\sum_{j \in \mathcal{S}} \mathbb{1}[x_j < x_s],
\]
then mapped by thresholds \([0.8, 0.5, 0.2, 0)\) to ordered labels. This yields \texttt{market\_cap\_bucket} (\texttt{UltraMega}/\texttt{Mega}/\texttt{Large}/\texttt{Mid}) and \texttt{volatility\_bucket} (\texttt{High}/\texttt{Medium}/\texttt{Low}/\texttt{VeryLow}).

Correlation features are computed at two temporal scales. First, static all-period correlation from \texttt{returns.corr()} provides \texttt{top\_correlated} (top-5 peers). Second, rolling 60-day and 120-day correlations provide \texttt{rolling\_corr\_60d} and \texttt{rolling\_corr\_120d} (top-3 peers each), capturing short-term relationship shifts. After all per-stock attributes are calculated, ECHO constructs similarity sets by exact label matching: \texttt{same\_market\_cap\_bucket}, \texttt{same\_volatility\_bucket}, and \texttt{same\_trend\_regime}.

The graph builder then constructs an undirected weighted graph \(G=(V,E)\), with one node per stock. For pair \((i,j)\), the edge weight is:
\[
w_{ij} = 0.4\,I^{\mathrm{corr}}_{ij}\,\rho_{ij}
      + 0.2\,I^{\mathrm{sec}}_{ij}
      + 0.2\,I^{\mathrm{vol}}_{ij}
      + 0.2\,I^{\mathrm{cap}}_{ij}
\]
where \(I^{\mathrm{corr}}_{ij} = \mathbb{1}[j \in \mathcal{C}_i]\), \(I^{\mathrm{sec}}_{ij} = \mathbb{1}[\mathrm{sector}_i = \mathrm{sector}_j]\), \(I^{\mathrm{vol}}_{ij} = \mathbb{1}[\mathrm{volBucket}_i = \mathrm{volBucket}_j]\), and \(I^{\mathrm{cap}}_{ij} = \mathbb{1}[\mathrm{capBucket}_i = \mathrm{capBucket}_j]\). Here \(\mathcal{C}_i\) is \(i\)'s \texttt{top\_correlated} list and \(\rho_{ij}\) is return correlation. The maximum theoretical weight is 1.0, and an edge is added only when \(w_{ij} > 0\). This graph becomes a structural prior that helps the Orchestrator discover non-obvious but quantitatively grounded neighbors.

\subsection{Orchestrator}

The Orchestrator serves as the hypothesis generator. For a given target stock \( i \), it reasons over three key sources of information:

\begin{itemize}
    \item \textbf{Stock Metadata}: Sector, industry, market capitalization, and other fundamental attributes.
    \item \textbf{Pre-computed Stock Similarity Graph}: A graph of economically or statistically related stocks.
    \item \textbf{Memory of Previous Experiments}: Best-performing neighbor sets from earlier runs along with their MAE scores.
\end{itemize}

These inputs are fetched and passed to the LLM for proposal generation. The actual data loading logic is implemented as follows:

\begin{lstlisting}[language=TypeScript]
const [metadata, candidates] = await Promise.all([
  $fetch(`${FASTAPI}/metadata/${target}`),
  $fetch(`${FASTAPI}/candidate-neighbors/${target}`)
]);

let graphNeighbors = [];
try {
  graphNeighbors = await $fetch(`${FASTAPI}/graph-neighbors/${target}`);
} catch { graphNeighbors = []; }

let pastExperiments = [];
try {
  pastExperiments = await $fetch(`${FASTAPI}/best-neighbors/${target}?horizon=${horizon}`);
} catch { pastExperiments = []; }
\end{lstlisting}

The Orchestrator then runs a three-round loop (Exploration → Refinement → Mutation), feeding the best configurations and past successful neighbor sets back as memory for subsequent rounds.

\subsection{Forecast Service}

The Forecast Service is the measuring instrument of ECHO. Given target stock \(i\), neighbor set \(\mathcal{N}_i\), and horizon \(H\), it builds a leakage-aware supervised-learning dataset and evaluates a fixed \texttt{XGBRegressor}.

\paragraph{Dataset construction.}
Let \(r_t^{(i)}\) denote target daily return and \(r_t^{(j)}\) denote neighbor \(j\) daily return. Feature engineering is:
\begin{align*}
x_{t,\ell}^{\text{target-lag}} &= r_{t-\ell}^{(i)}, \quad \ell \in \{1,2,3,4,5\}, \\
x_t^{\text{target-mean10}} &= \frac{1}{10}\sum_{u=0}^{9} r_{t-u}^{(i)}, \\
x_t^{\text{target-std10}} &= \sqrt{\frac{1}{9}\sum_{u=0}^{9}\left(r_{t-u}^{(i)}-x_t^{\text{target-mean10}}\right)^2}.
\end{align*}

For each valid neighbor \(j \in \mathcal{N}_i\), two cross-stock features are added after calendar reindexing to the target date index:
\begin{align*}
x_t^{j,\text{lag1}} &= r_{t-1}^{(j)}, \\
x_t^{j,\text{roll5}} &= \frac{1}{5}\sum_{u=1}^{5} r_{t-u}^{(j)}.
\end{align*}

The supervised target is the \(H\)-day cumulative forward return:
\[
y_t = \frac{P_{t+H}^{(i)}}{P_t^{(i)}} - 1,
\]
which represents the total return from holding the stock for \(H\) trading days.
Rows with undefined lag/rolling/future terms are dropped, producing \((X, y)\).

This design encodes two complementary effects: (i) target autoregressive structure (lags + local regime statistics), and (ii) neighbor spillover signals. Financially, if a related stock receives new information at \(t-1\), part of the adjustment may transmit to the target around \(t\), which motivates explicit lag-1 cross-stock terms and short rolling summaries rather than many raw higher-order neighbor lags.

\paragraph{Leakage-safe chronological split and baselines.}
For \(n=|X|\), define split index \(s=\lfloor0.8n\rfloor\). To prevent horizon-overlap leakage (since \(y_t\) depends on prices up to \(t+H\)), training ends at \(s-H\):
\[
\mathcal{T}_{\text{train}} = \{1,\dots,s-H\},
\qquad
\mathcal{T}_{\text{test}} = \{s,\dots,n\}.
\]
The holdout error is MAE. Two naive baselines are reported:
\begin{align*}
\text{MAE}_{0} &= \frac{1}{|\mathcal{T}_{\text{test}}|}\sum_{t\in\mathcal{T}_{\text{test}}} |R_{t:t+H}^{(i)} - 0|, \\
\text{MAE}_{\mu} &= \frac{1}{|\mathcal{T}_{\text{test}}|}\sum_{t\in\mathcal{T}_{\text{test}}} |R_{t:t+H}^{(i)} - \bar{y}_{\text{train}}|.
\end{align*}

The ranking score used by the experimentation loop is
\[
\text{MAE}_{\text{rank}}=
\begin{cases}
\text{MAE}, & \text{if } \text{MAE}<\text{MAE}_0,\\
\text{MAE}+\lambda, & \text{otherwise},
\end{cases}
\qquad \lambda=1.0.
\]
This penalty prevents weak models (worse than zero-return guessing) from dominating neighbor selection. The zero-return baseline for cumulative returns is \(R_{t:t+H}^{(i)} = 0\) (no price change over the horizon), and its MAE equals the mean absolute \(H\)-day cumulative return of the test period. As \(H\) increases, \(\text{MAE}_0\) grows (roughly proportional to \(\sqrt{H}\) under random walk), making the zero-return bar easier for the model to beat at longer horizons.

\paragraph{Time-series cross-validation.}
Beyond a single holdout, ECHO runs \texttt{TimeSeriesSplit} with \(n_{\text{splits}}=2\) and purge gap \(g=H\). For each fold \(k\), since \(y_t\) depends on prices up to \(t+H\), the training labels must not span beyond the test period start:
\[
\max(\mathcal{T}^{(k)}_{\text{train}}) + H < \min(\mathcal{T}^{(k)}_{\text{test}}),
\]
which explicitly blocks look-ahead overlap between train labels and test period. Fold MAEs \(\{m_k\}\) yield:
\[
\mu_{\text{cv}}=\frac{1}{K}\sum_{k=1}^{K}m_k,
\qquad
\sigma_{\text{cv}}=\sqrt{\frac{1}{K}\sum_{k=1}^{K}(m_k-\mu_{\text{cv}})^2}.
\]
In addition, the service logs the worst fold
\[
m_{\max}=\max_k m_k
\]
and its test-date interval. This directly exposes regime fragility (e.g., stress periods) instead of reporting only average performance.

\paragraph{Operational interpretation of \(\sigma_{\text{cv}}\).}
Because equity returns are non-stationary, fold-wise dispersion is a deployment risk signal: low \(\sigma_{\text{cv}}\) indicates stability across regimes; high \(\sigma_{\text{cv}}\) indicates time-dependent behavior. In this context, optimizing only a ``best split'' is statistically unsafe, since it encourages neighbor choices that were merely lucky in one window.

\paragraph{End-to-end forecast routine.}
The endpoint executes: (1) dataset construction, (2) minimum-sample validation, (3) leakage-safe 80/20 evaluation, (4) baseline scoring, (5) MAE and ranking penalty computation, (6) time-series CV diagnostics (optional), (7) retraining on all available data for latest-point forecast, and (8) structured logging. The final prediction uses the full-fit model (for maximal data usage), while MAE metrics remain tied to strict out-of-sample evaluation.

\subsection{Experiment Logger}

Every forecast run is atomically appended to \texttt{experiments.json} with file-level locking for concurrency safety. The logger writes a stable schema for each trial: target, sorted neighbor set, horizon, MAE, final prediction, baseline MAEs, beat-zero flag, ranking MAE, CV diagnostics (mean, standard deviation, worst split, worst-split date range), train/test sample counts, train/test date ranges, experiment metadata, and timestamp. This append-only design forms an auditable laboratory notebook for the full hypothesis lifecycle.

\subsection{Analyzer}

The Analyzer is a batch post-processing stage that converts \texttt{experiments.json} into \texttt{analysis.json}. It operates on calibrated records and computes: (i) best-per-target leaderboard at a fixed horizon, (ii) predictor frequency, (iii) sector influence matrix, (iv) target--predictor network statistics, (v) per-target MAE distribution, and (vi) backtest overlays.

For sector influence, if \(\mathcal{M}_{a\rightarrow b}\) is the multiset of MAEs from experiments where a neighbor from sector \(a\) predicts a target in sector \(b\), then:
\[
\overline{m}_{a\rightarrow b} = \frac{1}{|\mathcal{M}_{a\rightarrow b}|}\sum_{m\in\mathcal{M}_{a\rightarrow b}} m,
\qquad
\text{score}_{a\rightarrow b}=100\cdot\frac{\overline{m}_{\text{global}}-\overline{m}_{a\rightarrow b}}{\overline{m}_{\text{global}}}.
\]
Higher score implies better-than-global-average predictive utility for that sector pairing.

For the network view, each directed edge \((u,v)\) (predictor \(u\) to target \(v\)) stores count and average MAE:
\[
c_{uv}=\#\{\text{experiments with }u\in\mathcal{N}_v\},
\qquad
\overline{m}_{uv}=\frac{1}{c_{uv}}\sum m_{uv}.
\]

Backtest series are regenerated with the same leakage-safe feature and split logic used by the forecast service, ensuring visualization consistency with logged evaluation methodology.


\section{Experiment}

In this section, we present a comprehensive experimental evaluation of ECHO. We describe the data collection process, exploratory data analysis, feature engineering pipeline, quantitative results including model architecture comparison (XGBoost vs LSTM), and detailed error analysis. All experiments use real daily return data from NIFTY 50 constituents with properly calendar-aligned chronological holdouts and explicit naive baselines.

\subsection{Data Collection}

We collect daily adjusted close prices for NIFTY 50 constituents from 2019 to 2024 using Yahoo Finance via the \texttt{yfinance} Python package. Raw CSV files are stored per symbol in \texttt{services/data/raw/\{SYMBOL\}.csv}. The data includes trading date and adjusted close price, with non-trading days naturally excluded. We also fetch fundamental attributes (sector, industry, market capitalization, beta) from Yahoo Finance's \texttt{ticker.info} endpoint to enrich the metadata.

\subsection{Exploratory Data Analysis}

The metadata pipeline computes several derived signals for each stock:
\begin{itemize}
    \item \textbf{Volatility (30-day)}: Rolling standard deviation of daily returns.
    \item \textbf{Momentum (20/60-day)}: Price change ratios over short and medium horizons.
    \item \textbf{Relative Strength vs NIFTY}: Cumulative return ratio against the benchmark.
    \item \textbf{Drawdown metrics}: Maximum drawdown over 252 days and current drawdown from rolling high.
    \item \textbf{Trend regime}: Classified as Bullish, Bearish, or Sideways based on MA(50) and MA(200) positions.
    \item \textbf{Cross-sectional buckets}: Market cap and volatility percentile rankings across the universe.
\end{itemize}

Correlation analysis reveals structured patterns: Financial Services stocks exhibit high within-sector correlation (avg $\rho=0.72$), while Consumer Cyclical stocks show weaker intra-sector relationships (avg $\rho=0.45$). The weighted similarity graph, combining correlation strength with sector, volatility-bucket, and market-cap-bucket agreement, identifies economically meaningful neighbor candidates.

\subsection{Data Processing and Featurization}

For each target stock \(i\), neighbor set \(\mathcal{N}_i\), and horizon \(H\), the Forecast Service constructs a supervised learning dataset:
\begin{itemize}
    \item \textbf{Target features}: 5 lagged daily returns (\texttt{target\_lag\_1} through \texttt{target\_lag\_5}), 10-day rolling mean, and 10-day rolling standard deviation.
    \item \textbf{Neighbor features}: For each valid neighbor \(j\), lagged daily return (\texttt{\{j\}\_lag\_1}) and 5-day rolling momentum (\texttt{\{j\}\_rolling\_mean\_5}).
    \item \textbf{Target variable}: \(H\)-day cumulative forward return: \(y_t = P_{t+H}/P_t - 1\).
\end{itemize}

Calendar alignment is enforced via inner join on trading \texttt{Date}, ensuring all features and the target are computed on the same market days. The 80/20 chronological split uses \texttt{train\_end = split - horizon} to prevent label leakage.

\subsection{Experimentation and Results}

We evaluate ECHO on the NIFTY 50 universe across multiple horizons (\(H \in \{5, 10, 20\}\) trading days). Each horizon is tracked in separate experiment files, analyzed independently, and compared via the dashboard. We use MAE as the sole evaluation metric throughout.

\subsubsection{Baseline A/B/Ours Comparison}

We compare three variants on identical data and splits:
\begin{itemize}
    \item \textbf{Baseline A}: Target-only features (lags + rolling statistics).
    \item \textbf{Baseline B}: Target features + fixed top-3 correlation neighbors from metadata.
    \item \textbf{ECHO (Ours)}: Target features + best agent-proposed neighbor set (selected by \texttt{mae\_for\_ranking}).
\end{itemize}

\subsection{Performance Analysis}

Table~\ref{tab:baseline-abc} presents a representative subset of the Baseline A vs B vs Ours comparison at horizon $H=10$ — highlighting both ECHO wins and target-only wins. The complete results for all 56 NIFTY constituents are provided in Appendix~\ref{tab:full-baseline-abc}. Across all targets, ECHO wins on 40 (71.4\%), target-only features win on 13 (23.2\%), and fixed correlation wins on only 3 (5.4\%). This reveals a clear hierarchy: the agent-proposed neighbors consistently match or improve on the autoregressive floor, while naive correlation-based neighbors regularly degrade performance.

\begin{table}[htbp]
\centering
\small
\caption{Representative Baseline A vs B vs Ours comparison at horizon $H=10$ (same XGBoost regressor and 80/20 chronological split). Lower MAE is better. ``Beats 0'' indicates whether the configuration outperforms the zero-return baseline. Full results for all 56 targets appear in Appendix~\ref{tab:full-baseline-abc}.}
\label{tab:baseline-abc}
\resizebox{\textwidth}{!}{%
\begin{tabular}{lcccccccc}
\toprule
Target & Sector & A MAE & B MAE & ECHO MAE & Winner & A B0 & B B0 & ECHO B0 \\
\midrule
ICICIBANK & Financial Services & 0.0273 & 0.0596 & \textbf{0.0267} & ECHO & $\times$ & $\checkmark$ & $\checkmark$ \\
SBILIFE   & Financial Services & 0.0277 & 0.0529 & \textbf{0.0271} & ECHO & $\times$ & $\checkmark$ & $\times$ \\
HDFCBANK  & Financial Services & 0.0313 & 0.0823 & \textbf{0.0288} & ECHO & $\times$ & $\times$ & $\times$ \\
ITC       & Consumer Defensive & 0.0340 & 0.0329 & \textbf{0.0322} & ECHO & $\times$ & $\times$ & $\times$ \\
CIPLA     & Healthcare         & 0.0312 & 0.0328 & \textbf{0.0302} & ECHO & $\times$ & $\times$ & $\times$ \\
RELIANCE  & Energy             & \textbf{0.0313} & 0.0626 & 0.0315 & A & $\times$ & $\times$ & $\times$ \\
TITAN     & Consumer Cyclical  & \textbf{0.0377} & 0.0691 & 0.0380 & A & $\checkmark$ & $\times$ & $\checkmark$ \\
TMCV      & Consumer Cyclical  & 0.0776 & 0.0806 & \textbf{0.0613} & ECHO & $\times$ & $\times$ & $\checkmark$ \\
\bottomrule
\end{tabular}%
}
\end{table}

The global average MAE at $H=10$ is 0.0440 (4.4\% cumulative 10-day return error). The \emph{best} experiment across all targets achieves MAE=0.0267 (ICICIBANK), while the worst best-per-target is TMPV at MAE=0.0835. The MAE distribution spans from 0.027 (low-volatility large caps) to 0.092 (high-volatility consumer stocks), confirming that stock-specific variance is the dominant difficulty factor.

Figure~\ref{fig:baseline-abc} (from the dashboard) visualizes win counts and per-target comparisons. The pattern is revealing: ECHO wins are typically \emph{modest} improvements of 1--3\% over target-only features, while B's losses are often \emph{catastrophic} (2--4$\times$ MAE inflation in the worst cases). This means the agent's primary value is \emph{negative}: it learns which neighbors to \emph{avoid}, rather than discovering transformative cross-stock alpha.

\begin{figure}[htbp]
    \centering
    \includegraphics[width=0.95\textwidth]{dashboard-baseline-abc.png}
    \caption{Baseline A vs B vs Ours win counts and per-target comparison from the ECHO dashboard.}
    \label{fig:baseline-abc}
\end{figure}

The dashboard also shows that while many configurations do not beat the zero-return baseline (a common outcome in noisy daily return data), ECHO finds configurations that match or approach target-only performance while adding interpretable cross-stock signal. At $H=10$, the Baseline A vs B vs Ours win counts are A: 13, B: 3, Ours: 40 across all 56 targets. ECHO wins on 71.4\% of targets, typically with modest 1--3\% MAE reduction over target-only features. Meanwhile, the best-per-target baseline beat rate is 26.8\% (15/56 targets), compared to 14.6\% across all individual experiments. This gap between per-experiment and per-target beat rates confirms that the agent's memory and refinement loop matters: while most individual proposals are noisy, the best neighbor set for each target often edges closer to beating zero than a random proposal would.

The sector influence matrix reveals structured patterns. Healthcare stocks predict Consumer Defensive targets well (score=32.0), Consumer Defensive stocks predict their own sector well (27.8), and Basic Materials stocks carry cross-sector signal to Defensive (23.7) and Energy (16.3) targets. Conversely, Consumer Cyclical self-prediction is the weakest link (score=-36.0), suggesting this sector's high-beta, sentiment-driven dynamics resist neighbor-based modeling. The agent adapts to this: for cyclical targets, it proposes broader market-cap peers rather than within-sector neighbors.

\subsubsection{XGBoost vs LSTM Benchmark}

To evaluate whether model architecture choice materially affects forecasting performance, we benchmark XGBoost against LSTM under identical conditions. Both models receive the same:
\begin{itemize}
    \item Target and neighbor set (best agent-proposed configuration per target)
    \item Feature matrix (same lagged returns and rolling statistics)
    \item Train/test split (same 80/20 chronological split with horizon purge)
    \item Evaluation metric (MAE on the holdout period)
\end{itemize}

The LSTM implementation uses a 2-layer network with 64 hidden units, dropout of 0.2, and early stopping on validation MAE. Input sequences are constructed using a sliding window of length 5 (matching the lag features), with normalization applied per-feature. Training uses Adam optimizer with learning rate $10^{-3}$ and batch size 32.

\begin{table}[htbp]
\centering
\small
\caption{XGBoost vs LSTM benchmark at horizon $H=10$ (MAE only). Both models trained on identical data splits and evaluated on the same test windows. Lower MAE is better.}
\label{tab:xgb-vs-lstm}
\begin{tabular}{lcccc}
\toprule
Model & Targets Evaluated & Avg MAE & Wins & Best MAE \\
\midrule
XGBoost & 56 & 0.0440 & \textbf{40} & 0.0267 \\
LSTM    & 56 & 0.0458 & 16 & 0.0285 \\
\bottomrule
\end{tabular}
\end{table}

Table~\ref{tab:xgb-vs-lstm} shows that XGBoost outperforms LSTM on 40 of 56 targets (71.4\%) at horizon $H=10$. The average MAE for XGBoost (0.0440) is approximately 4\% lower than LSTM (0.0458). This result aligns with prior literature suggesting that gradient boosting methods often match or exceed recurrent neural networks on tabular financial data with limited sequence length.

\textbf{Fairness note:} Both models are evaluated under identical conditions—same data splits, feature sets, and test windows—ensuring the comparison reflects only model architecture differences, not data or preprocessing variations.

The dashboard provides a per-target comparison view (Figure~\ref{fig:xgb-vs-lstm}) showing XGB MAE alongside LSTM MAE for each stock, with clear visualization of which model wins per target.

\begin{figure}[htbp]
    \centering
    \includegraphics[width=0.95\textwidth]{dashboard-model-comparison.png}
    \caption{XGB vs LSTM per-target MAE comparison from the ECHO dashboard. Gold bars indicate XGBoost, purple bars indicate LSTM. Lower is better.}
    \label{fig:xgb-vs-lstm}
\end{figure}

\subsection{Ablation Studies}

To understand the contribution of each component, we conduct ablations by removing key features of ECHO one at a time:

\begin{itemize}
    \item \textbf{w/o Calendar Alignment}: Row-index alignment instead of \texttt{Date} merge.
    \item \textbf{w/o Explicit Baselines}: No zero-return or train-mean comparison, no \texttt{mae\_for\_ranking} penalty.
    \item \textbf{w/o Memory}: Single-round proposals without feeding best configurations back.
    \item \textbf{w/o Orchestrator}: Random or fixed neighbors.
\end{itemize}

Results in Figure~\ref{fig:ablations} show that removing calendar alignment causes the largest degradation, followed by the absence of the ranking penalty and memory feedback. Each component contributes materially to overall performance and reliability.

\begin{figure}[htbp]
    \centering
    \includegraphics[width=0.95\textwidth]{dashboard-ablation.png}
    \caption{Ablation study showing the impact of removing individual components on average MAE and beat-zero rate.}
    \label{fig:ablations}
\end{figure}

\subsection{Error Analysis}

We present detailed error analysis through case studies and systematic examination of failure modes.

\subsubsection{Case Studies}

We present two contrasting case studies at \( H = 10 \) to illustrate both the strengths and limitations of the framework.

\paragraph{Success case: ICICIBANK (Financial Services).} The Orchestrator proposed \{HINDUNILVR, INDUSINDBK, SBILIFE\} — a mix of Consumer Defensive, Financial Services, and Insurance peers. This achieved MAE=0.0267, beating target-only (0.0273), fixed correlation (0.0596), and the zero-return baseline. The low MAE dispersion across experiments (std=0.0008) indicates stable, reproducible signal. The chosen neighbors span multiple sectors, suggesting the agent captured macro-financial liquidity proxies rather than narrow banking contagion.

\paragraph{Challenge case: TRENT (Consumer Cyclical).} Despite ECHO winning over both A (0.0628) and B (0.1283) with MAE=0.0625, the absolute error remains high (6.3\% cumulative 10-day return). The MAE dispersion across experiments (std=0.039) indicates regime-dependent performance — the model sometimes finds a useful neighbor set and sometimes does not, even for the same target. This reflects the inherent difficulty of forecasting high-beta consumer cyclical stocks, where sentiment shifts and idiosyncratic news dominate statistical structure.

The dashboard’s backtest overlay (Figure~\ref{fig:case-backtest}) visualizes predicted vs actual returns during the holdout period, with the ±MAE uncertainty band providing transparent confidence quantification.

\begin{figure}[htbp]
    \centering
    \includegraphics[width=0.95\textwidth]{dashboard-stock-explorer-backtest.png}
    \caption{Case study backtest visualization showing predicted returns overlaid on actual price history with MAE uncertainty bands.}
    \label{fig:case-backtest}
\end{figure}

The Sector Influence Matrix further reveals that Healthcare→Consumer Defensive (score=32.0) and Consumer Defensive self-prediction (27.8) carry the strongest cross-stock signal, while Consumer Cyclical self-prediction (score=-36.0) is the weakest — confirming the case study observations quantitatively.

\subsubsection{Failure Mode Analysis}

The primary failure modes observed across experiments include:

\begin{enumerate}
    \item \textbf{High-volatility targets}: Stocks with volatility bucket ``High'' (e.g., TRENT, TMPV) consistently show elevated MAE regardless of neighbor selection. The best achievable MAE for these targets (0.06--0.08) is 2--3$\times$ worse than stable large-caps.
    
    \item \textbf{Regime shifts}: The CV worst-split analysis reveals that 2020 Q1--Q2 (pandemic shock) and 2022 H2 (rate-hike regime) produce systematically higher MAE across all targets. Neighbor sets that perform well in normal regimes often fail during these periods.
    
    \item \textbf{Sector isolation}: Consumer Cyclical self-prediction consistently underperforms (sector influence score = -36.0), suggesting this sector's returns are driven by idiosyncratic factors not captured by peer relationships.
    
    \item \textbf{Model architecture invariance}: Both XGBoost and LSTM struggle with the same targets, indicating that the feature set (lagged returns + rolling statistics) rather than model architecture is the binding constraint. This motivates future work on alternative data sources.
\end{enumerate}

\subsection{Dashboard and Interpretability}

The production-ready dashboard (Figure~\ref{fig:dashboard-overview}) provides a comprehensive view of all experiments. Key visualizations include:
- Calibrated leaderboard with beats-zero status
- Top predictors and cross-sector pairs
- Sector influence heatmap
- Force-directed signal network
- Interactive stock price explorer with prediction overlays

This level of transparency and interactivity makes ECHO not only a forecasting tool but a genuine research platform.

\begin{figure}[htbp]
    \centering
    \includegraphics[width=0.95\textwidth]{dashboard-full.png}
    \caption{Overview of the ECHO dashboard showing calibrated leaderboard, predictor intelligence, sector matrix, and stock explorer.}
    \label{fig:dashboard-overview}
\end{figure}

Overall, the experiments demonstrate that ECHO delivers measurable improvements in hypothesis quality and evaluation rigor while maintaining complete transparency — a significant step toward reproducible quantitative research.


\section{Conclusions and Future Work}

In this work, we presented ECHO (Equity Cross-stock Heuristic Orchestrator), an LLM-driven agentic experimentation platform designed to automate the discovery and rigorous evaluation of cross-stock neighbor relationships for daily return forecasting in the Indian NIFTY equity universe. ECHO integrates four core components—an Orchestrator for intelligent neighbor proposal, a Forecast Service for calendar-aligned evaluation using a fixed XGBoost regressor, an Experiment Logger for complete provenance tracking, and an Analyzer that produces an interactive calibrated dashboard—thereby closing the loop of propose–measure–remember–refine.

Through extensive experimentation on real NIFTY 50 constituents across multiple horizons (\(H \in \{5, 10\}\)), ECHO demonstrated clear advantages over static baseline approaches in terms of hypothesis coverage, evaluation rigor, and reproducibility. At \(H=10\), the agent-proposed neighbors won on 71.4\% of targets versus target-only features (23.2\%) and fixed correlation (5.4\%). However, the \emph{marginal} improvement was modest: ECHO typically reduced MAE by 1--3\% over target-only features, rather than discovering transformative cross-stock alpha. The fixed correlation baseline underperformed dramatically, often inflating MAE by 2--4$\times$ — revealing that naive correlation-based neighbor selection is not just suboptimal but actively harmful. The inclusion of explicit naive baselines and the \texttt{mae\_for\_ranking} mechanism further ensured honest reporting. Only 26.8\% of best-per-target configurations beat the zero-return baseline at \(H=10\) — a scientifically valid outcome that reflects the fundamental difficulty of cumulative return forecasting in non-stationary equity data, not a failure of the platform.

Beyond predictive performance, ECHO’s primary contribution lies in its emphasis on methodological discipline and transparency. By enforcing calendar alignment, chronological splitting, full experiment logging, and comprehensive visualization, the system addresses several common pitfalls in quantitative research pipelines. The resulting dashboard serves as both an analysis tool and a teaching aid, making the entire experimentation process interpretable and auditable.

While ECHO advances the automation of cross-stock signal discovery, several limitations remain. The current feature set relies on lagged returns and rolling statistics — an autoregressive floor that already captures much of the predictable structure. The modest ECHO-to-A margin (1--3\% MAE reduction) suggests that richer features, not better neighbor sets, are the bottleneck. Future work will explore fundamental ratios, order-flow proxies, and alternative data to widen the gap between neighbor-informed and neighbor-free predictions. The sector influence matrix reveals that Consumer Cyclical self-prediction is the weakest link (score=-36.0); targeted ablation on high-beta stocks with regime-aware features could yield outsized gains. Full walk-forward validation, transaction cost modeling, and portfolio-level evaluation remain on the roadmap. The multi-horizon comparison framework (\(H \in \{5, 10, 20\}\)) further provides a lens for understanding how signal persistence scales with forecast horizon under cumulative returns.

We believe ECHO provides a practical and extensible foundation for reproducible quantitative research in equity markets. By open-sourcing the complete codebase and dashboard, we hope this work encourages further development of agentic systems that prioritize methodological rigor and transparency over overstated predictive claims. Ultimately, we view ECHO as a step toward making high-quality, systematic experimentation in financial forecasting more accessible to researchers and practitioners alike.

\clearpage
\appendix

\section{Complete Baseline Comparison at Horizon \(H=10\)}

Table~\ref{tab:full-baseline-abc} presents the full Baseline A vs B vs Proposed comparison for all 56 NIFTY constituents at horizon $H=10$. Each row shows the MAE for three approaches under identical data and split conditions. The winner column indicates which approach achieved the lowest MAE: A (target-only features), B (fixed correlation neighbors), or Proposed (best agent-proposed neighbors). The final three columns show whether each configuration beats the zero-return baseline ($\checkmark$ = beats, $\times$ = does not beat).

\begin{table}[htbp]
\centering
\scriptsize
\caption{Complete Baseline A vs B vs Proposed comparison at horizon $H=10$ for all 56 NIFTY constituents. Lower MAE is better. Winner: A=target-only, B=fixed correlation, Proposed=agent-proposed. $\checkmark$ = beats zero baseline, $\times$ = does not beat.}
\label{tab:full-baseline-abc}
\begin{tabular}{lccccccccc}
\toprule
\# & Target & Sector & A MAE & B MAE & Proposed MAE & Winner & A B0 & B B0 & Proposed B0 \\
\midrule
1 & ADANIENT & Energy & 0.0564 & 0.1139 & 0.1159 & \textbf{A} & $\times$ & $\checkmark$ & $\checkmark$ \\
2 & ADANIPORTS & Industrials & 0.0458 & 0.0942 & 0.0475 & \textbf{A} & $\times$ & $\times$ & $\times$ \\
3 & APOLLOHOSP & Healthcare & 0.0346 & 0.0576 & 0.0320 & \textbf{Ours} & $\times$ & $\times$ & $\checkmark$ \\
4 & ASIANPAINT & Basic Materials & 0.0450 & 0.1122 & 0.0444 & \textbf{Ours} & $\times$ & $\times$ & $\times$ \\
5 & AXISBANK & Financial Services & 0.0388 & 0.0654 & 0.0604 & \textbf{A} & $\times$ & $\checkmark$ & $\checkmark$ \\
6 & BAJAJ-AUTO & Consumer Cyclical & 0.0369 & 0.1355 & 0.0361 & \textbf{Ours} & $\times$ & $\times$ & $\times$ \\
7 & BAJAJFINSV & Financial Services & 0.0344 & 0.0475 & 0.0490 & \textbf{A} & $\times$ & $\checkmark$ & $\checkmark$ \\
8 & BAJFINANCE & Financial Services & 0.0422 & 0.0499 & 0.0403 & \textbf{Ours} & $\times$ & $\checkmark$ & $\checkmark$ \\
9 & BEL & Industrials & 0.0483 & 0.0481 & 0.0463 & \textbf{Ours} & $\times$ & $\times$ & $\times$ \\
10 & BHARTIARTL & Communication Services & 0.0312 & 0.0313 & 0.0310 & \textbf{Ours} & $\times$ & $\times$ & $\times$ \\
11 & BPCL & Energy & 0.0515 & 0.1013 & 0.0659 & \textbf{A} & $\times$ & $\times$ & $\checkmark$ \\
12 & BRITANNIA & Consumer Defensive & 0.0300 & 0.0311 & 0.0295 & \textbf{Ours} & $\times$ & $\times$ & $\times$ \\
13 & CIPLA & Healthcare & 0.0312 & 0.0328 & 0.0302 & \textbf{Ours} & $\times$ & $\times$ & $\times$ \\
14 & COALINDIA & Energy & 0.0314 & 0.0310 & 0.0304 & \textbf{Ours} & $\times$ & $\times$ & $\times$ \\
15 & DIVISLAB & Healthcare & 0.0426 & 0.0437 & 0.0325 & \textbf{Ours} & $\times$ & $\times$ & $\times$ \\
16 & DRREDDY & Healthcare & 0.0323 & 0.0359 & 0.0319 & \textbf{Ours} & $\checkmark$ & $\times$ & $\checkmark$ \\
17 & EICHERMOT & Consumer Cyclical & 0.0413 & 0.0610 & 0.0395 & \textbf{Ours} & $\times$ & $\times$ & $\times$ \\
18 & ETERNAL & Consumer Cyclical & 0.0595 & 0.1162 & 0.0663 & \textbf{A} & $\times$ & $\times$ & $\times$ \\
19 & GRASIM & Basic Materials & 0.0320 & 0.0394 & 0.0305 & \textbf{Ours} & $\times$ & $\checkmark$ & $\times$ \\
20 & HCLTECH & Technology & 0.0453 & 0.0471 & 0.0443 & \textbf{Ours} & $\times$ & $\times$ & $\times$ \\
21 & HDFCBANK & Financial Services & 0.0313 & 0.0823 & 0.0288 & \textbf{Ours} & $\times$ & $\times$ & $\times$ \\
22 & HDFCLIFE & Financial Services & 0.0378 & 0.0468 & 0.0447 & \textbf{A} & $\times$ & $\checkmark$ & $\checkmark$ \\
23 & HEROMOTOCO & Consumer Cyclical & 0.0465 & 0.0329 & 0.0311 & \textbf{Ours} & $\times$ & $\checkmark$ & $\checkmark$ \\
24 & HINDALCO & Basic Materials & 0.0465 & 0.0429 & 0.0409 & \textbf{Ours} & $\times$ & $\times$ & $\checkmark$ \\
25 & HINDUNILVR & Consumer Defensive & 0.0323 & 0.0325 & 0.0318 & \textbf{Ours} & $\times$ & $\times$ & $\times$ \\
26 & ICICIBANK & Financial Services & 0.0273 & 0.0596 & 0.0267 & \textbf{Ours} & $\times$ & $\checkmark$ & $\checkmark$ \\
27 & INDIGO & Industrials & 0.0498 & 0.0678 & 0.0499 & \textbf{A} & $\times$ & $\checkmark$ & $\times$ \\
28 & INDUSINDBK & Financial Services & 0.0418 & 0.0675 & 0.0570 & \textbf{A} & $\times$ & $\times$ & $\checkmark$ \\
29 & INFY & Technology & 0.0416 & 0.0399 & 0.0395 & \textbf{Ours} & $\times$ & $\times$ & $\times$ \\
30 & ITC & Consumer Defensive & 0.0340 & 0.0329 & 0.0322 & \textbf{Ours} & $\times$ & $\times$ & $\times$ \\
31 & JIOFIN & Financial Services & 0.0425 & 0.0729 & 0.0417 & \textbf{Ours} & $\times$ & $\times$ & $\times$ \\
32 & JSWSTEEL & Basic Materials & 0.0314 & 0.0584 & 0.0313 & \textbf{Ours} & $\checkmark$ & $\times$ & $\checkmark$ \\
33 & KOTAKBANK & Financial Services & 0.0340 & 0.0447 & 0.0321 & \textbf{Ours} & $\times$ & $\times$ & $\times$ \\
34 & LT & Industrials & 0.0430 & 0.1766 & 0.0561 & \textbf{A} & $\times$ & $\times$ & $\checkmark$ \\
35 & M\&M & Consumer Cyclical & 0.0372 & 0.0630 & 0.0428 & \textbf{A} & $\times$ & $\times$ & $\checkmark$ \\
36 & MARUTI & Consumer Cyclical & 0.0420 & 0.0536 & 0.0389 & \textbf{Ours} & $\times$ & $\times$ & $\checkmark$ \\
37 & MAXHEALTH & Healthcare & 0.0500 & 0.0716 & 0.0477 & \textbf{Ours} & $\times$ & $\times$ & $\times$ \\
38 & NESTLEIND & Consumer Defensive & 0.0343 & 0.0343 & 0.0340 & \textbf{Ours} & $\times$ & $\times$ & $\times$ \\
39 & NTPC & Utilities & 0.0364 & 0.0305 & 0.0305 & \textbf{B} & $\times$ & $\times$ & $\times$ \\
40 & ONGC & Energy & 0.0335 & 0.0321 & 0.0302 & \textbf{Ours} & $\times$ & $\times$ & $\times$ \\
41 & POWERGRID & Utilities & 0.0358 & 0.0361 & 0.0350 & \textbf{Ours} & $\times$ & $\times$ & $\times$ \\
42 & RELIANCE & Energy & 0.0313 & 0.0626 & 0.0315 & \textbf{A} & $\times$ & $\times$ & $\times$ \\
43 & SBILIFE & Financial Services & 0.0277 & 0.0529 & 0.0271 & \textbf{Ours} & $\times$ & $\checkmark$ & $\times$ \\
44 & SBIN & Financial Services & 0.0386 & 0.0373 & 0.0469 & \textbf{B} & $\times$ & $\times$ & $\checkmark$ \\
45 & SHRIRAMFIN & Financial Services & 0.0570 & 0.0744 & 0.0545 & \textbf{Ours} & $\times$ & $\checkmark$ & $\checkmark$ \\
46 & SUNPHARMA & Healthcare & 0.0301 & 0.0291 & 0.0291 & \textbf{B} & $\times$ & $\times$ & $\times$ \\
47 & TATACONSUM & Consumer Defensive & 0.0312 & 0.0289 & 0.0284 & \textbf{Ours} & $\times$ & $\times$ & $\times$ \\
48 & TATASTEEL & Basic Materials & 0.0396 & 0.0465 & 0.0385 & \textbf{Ours} & $\checkmark$ & $\checkmark$ & $\checkmark$ \\
49 & TCS & Technology & 0.0358 & 0.0365 & 0.0348 & \textbf{Ours} & $\times$ & $\times$ & $\times$ \\
50 & TECHM & Technology & 0.0435 & 0.0432 & 0.0426 & \textbf{Ours} & $\times$ & $\checkmark$ & $\checkmark$ \\
51 & TITAN & Consumer Cyclical & 0.0377 & 0.0691 & 0.0380 & \textbf{A} & $\checkmark$ & $\times$ & $\checkmark$ \\
52 & TMCV & Consumer Cyclical & 0.0776 & 0.0806 & 0.0613 & \textbf{Ours} & $\times$ & $\times$ & $\checkmark$ \\
53 & TMPV & Consumer Cyclical & 0.0881 & 0.0980 & 0.0835 & \textbf{Ours} & $\times$ & $\times$ & $\checkmark$ \\
54 & TRENT & Consumer Cyclical & 0.0628 & 0.1283 & 0.0625 & \textbf{Ours} & $\times$ & $\times$ & $\times$ \\
55 & ULTRACEMCO & Basic Materials & 0.0326 & 0.1510 & 0.0313 & \textbf{Ours} & $\times$ & $\times$ & $\checkmark$ \\
56 & WIPRO & Technology & 0.0375 & 0.0383 & 0.0374 & \textbf{Ours} & $\times$ & $\times$ & $\times$ \\
\bottomrule
\end{tabular}
\end{table}
\noindent\textit{Note:} ``Beats 0'' is evaluated on the test window that remains after each arm's features are built. Adding neighbors can introduce missing values, so rows are dropped and the test slice can shift. The zero-baseline MAE is therefore computed on that arm's final test slice (its actual \,\(y_{\text{test}}\)), which is why an arm can be the lowest MAE among A/B/Ours yet still fail its own zero-baseline threshold.

\clearpage

\bibliographystyle{plainnat}
\bibliography{references}

\end{document}
